%-------------------------------------------------------------------------------
% Mattia Giovannini — English CV (public-safe), Awesome-CV layout
% Build with LuaLaTeX:  latexmk -lualatex CV_EN-public.tex
%-------------------------------------------------------------------------------
\documentclass[11pt, a4paper]{../awesome-cv}
% Shared Awesome-CV settings for the public CVs (IT/EN/FR).
% Each CV loads this after \documentclass{../awesome-cv}; build with LuaLaTeX.

\geometry{left=1.4cm, top=.9cm, right=1.4cm, bottom=1.6cm, footskip=.5cm}

% Same blue accent as the other CVs
\definecolor{awesome}{RGB}{80, 100, 200}
\setbool{acvSectionColorHighlight}{true}
\renewcommand{\acvHeaderSocialSep}{\quad\textbar\quad}

% No hyphenation: words stay whole when application portals extract the text
\hyphenpenalty=10000
\exhyphenpenalty=10000

\name{Mattia}{Giovannini}
\email{mattia.giovannini6@studio.unibo.it}
\homepage{mattia-giovannini.github.io}
\github{frigori}

% Re-load the class fonts with LuaTeX's node renderer: with HarfBuzz, explicit
% hyphens ("eye-tracking") are extracted from the PDF as invisible soft hyphens.
\defaultfontfeatures{Ligatures=TeX, Renderer=Node}
\setmainfont{Source Sans 3}[
  UprightFont=*, ItalicFont=* Italic, BoldFont=* Bold, BoldItalicFont=* Bold Italic,
  FontFace={l}{n}{Font=* Light}, FontFace={l}{it}{Font=* Light Italic},
]
\setsansfont{Source Sans 3}[
  UprightFont=*, ItalicFont=* Italic, BoldFont=* Bold, BoldItalicFont=* Bold Italic,
  FontFace={l}{n}{Font=* Light}, FontFace={l}{it}{Font=* Light Italic},
]
\renewfontfamily\roboto{Roboto}[
  UprightFont=*, ItalicFont=* Italic, BoldFont=* Bold, BoldItalicFont=* Bold Italic,
  FontFace={l}{n}{Font=* Light}, FontFace={l}{it}{Font=* Light Italic},
]


\position{MSc Student in Computer Engineering (Double Degree){\enskip\cdotp\enskip}BSc in Computer Engineering}
\address{Bologna, Italy}

\begin{document}

\makecvheader[C]
\makecvfooter{}{Mattia Giovannini~~~·~~~Curriculum Vitae}{\thepage}

%-------------------------------------------------------------------------------
\cvsection{Profile}

\begin{cvparagraph}
Computer Engineering graduate from the University of Bologna with a strong focus on artificial intelligence,
machine learning and data analysis, enrolled in a double-degree MSc in Computer Engineering (University of
Bologna) and Ubiquitous Systems and Networks (Université Côte d'Azur, Polytech Nice Sophia). My undergraduate
studies gave me solid foundations in algorithms, statistics and software engineering, which I applied in a
research context during an Erasmus semester in France. I am motivated by the challenge of designing
intelligent systems to solve real-world problems.
\end{cvparagraph}

%-------------------------------------------------------------------------------
\cvsection{Education}
\begin{cventries}
  \cventry
    {Double degree with the MSc in Ubiquitous Systems and Networks, Université Côte d'Azur}
    {MSc in Computer Engineering — Double Degree}
    {Bologna, Italy \& Nice, France}
    {2026 -- 2028 (Ongoing)}
    {}

  \cventry
    {Final grade: 94/110}
    {BSc in Computer Engineering — University of Bologna}
    {Bologna, Italy}
    {2021 -- July 2026}
    {
      \begin{cvitems}
        \item {\textbf{Thesis:} \textit{Autism Spectrum Disorder Detection via Eye Tracking Data using Machine Learning}, applying classification and signal-analysis techniques to clinical eye-tracking recordings for ASD screening support.}
        \item {Relevant courses: Mathematical Analysis, Linear Algebra \& Geometry, Probability \& Statistics, Fundamentals of Computer Science, Operating Systems, Computer Networks, Information Systems, Software Engineering, Web Technologies, Computer Architecture.}
      \end{cvitems}
    }

  \cventry
    {Erasmus study project}
    {Université de Tours}
    {Tours, France}
    {2025}
    {
      \begin{cvitems}
        \item {Developed a project on \textbf{machine learning for autism spectrum disorder (ASD) analysis}, applying classification and feature extraction to eye-tracking data.}
        \item {Worked in an international, French-speaking academic environment; strengthened data analysis, Python and scientific methodology.}
      \end{cvitems}
    }

  \cventry
    {Scientific High School Diploma}
    {Liceo Augusto Righi}
    {Bologna, Italy}
    {2016 -- 2021}
    {}
\end{cventries}

%-------------------------------------------------------------------------------
\cvsection{Projects}
\begin{cventries}
  \cventry
    {GenAI Project Team}
    {BOOM GEN AI Challenge 2026 — VEM Sistemi}
    {}
    {2026}
    {
      \begin{cvitems}
        \item {Collaborated on a \textbf{GenAI-based assistant} for automated tender analysis, supporting the \textbf{GO/NO-GO} evaluation process.}
        \item {Designed and implemented a \textbf{Python scoring module} to assess alignment between tender requirements and company capabilities.}
        \item {Used \textbf{prompt engineering} and \textbf{Azure OpenAI} for requirement extraction, classification and enrichment; analysed a document parsing pipeline for unstructured tenders.}
      \end{cvitems}
    }

  \cventry
    {Project work · Enterprise Platforms and DevOps M, University of Bologna}
    {Resource-Aware Kubernetes Scheduler}
    {}
    {2026 (Ongoing)}
    {
      \begin{cvitems}
        \item {Built a \textbf{custom Kubernetes scheduler} in Python that places each pod on the lowest-latency node with enough free CPU, memory and pod capacity.}
        \item {Measures node latency live, binds pods through the Kubernetes API and explains every decision with \textit{Scheduled}/\textit{FailedScheduling} events.}
        \item {Now extending it to \textbf{smart placement and replacement} of services on mobile edge nodes (humanoid robots moving between workstations).}
      \end{cvitems}
    }

  \cventry
    {Discrete and Network Optimization Problems M, University of Bologna}
    {SDN Link-Failure Recovery Benchmark}
    {\href{https://github.com/frigori/net-restoration-bench}{\faGithub\ net-restoration-bench}}
    {2026 (Ongoing)}
    {
      \begin{cvitems}
        \item {Compares \textbf{shared backup path protection (SBPP)}, \textbf{Fast Reroute} and \textbf{global MILP recomputation} on capacity overhead, restoration rate, reaction time and path stretch.}
        \item {Python with NetworkX and PuLP; one shared plan/react interface runs all strategies on identical topologies and failure scenarios.}
      \end{cvitems}
    }

  \cventry
    {Lean Startup Lab · Technology Entrepreneurship}
    {AgenticOS — Multi-Agent AI Audit}
    {\href{https://agenticos-mvp.netlify.app}{\faGlobe\ live MVP}}
    {2026}
    {
      \begin{cvitems}
        \item {Built a \textbf{Red Team / Blue Team multi-agent system} (Gemini) that audits why AI shopping agents skip a fashion product.}
        \item {A Blue Team agent compiles the fixes into \textbf{JSON-LD} and \textbf{llms.txt} and judges them in a sandbox; shipped as an MVP with web UI and containerised engine.}
      \end{cvitems}
    }

  \cventry
    {Personal project}
    {Self-Hosted Homelab}
    {\href{https://mattia-giovannini.github.io/homelab/}{\faGlobe\ homelab page}}
    {Ongoing}
    {
      \begin{cvitems}
        \item {Run containerised services on Raspberry Pi nodes with \textbf{Docker, rootless Podman and Kubernetes}, behind a TLS reverse proxy.}
        \item {Prometheus/Grafana monitoring, CrowdSec edge security; a testbed for giving agentic AI tools \textbf{scoped, auditable access}.}
      \end{cvitems}
    }
\end{cventries}

%-------------------------------------------------------------------------------
\cvsection{Skills}
\begin{cvskills}
  \cvskill{Programming}{C, Java, Python, JavaScript, HTML/CSS, SQL, Bash}
  \cvskill{ML \& Data}{scikit-learn, NumPy, pandas, Matplotlib, Jupyter Notebook}
  \cvskill{Systems \& Networks}{Linux, TCP/IP, socket programming, relational databases}
  \cvskill{Tools}{Git, \LaTeX, MATLAB, Docker, Podman, Kubernetes}
  \cvskill{Interests}{Machine learning, data analysis, containerization, system automation}
\end{cvskills}

%-------------------------------------------------------------------------------
\cvsection{Languages}
\begin{cvskills}
  \cvskill{Italian}{Native}
  \cvskill{English — C1}{Duolingo English Test 130/160 (C1) · CLA C1, University of Bologna · Cambridge First Certificate B2 (178/190)}
  \cvskill{French — C1}{CLA C1, University of Bologna · B2/C1 courses at Université de Tours}
\end{cvskills}

%-------------------------------------------------------------------------------
\cvsection{Other Experience}
\begin{cventries}
  \cventry
    {Student Webmaster \& Editor}
    {Liceo Augusto Righi}
    {Bologna, Italy}
    {2020 -- 2021}
    {
      \begin{cvitems}
        \item {Designed and managed the school's student platforms (\textit{Righi Network} and \textit{Il Sottobanco}), supporting the school newspaper's online presence.}
      \end{cvitems}
    }

  \cventry
    {Volunteer Summer Camp Educator}
    {Estate Ragazzi}
    {Bologna, Italy}
    {2016 -- 2020}
    {
      \begin{cvitems}
        \item {Group management, teamwork, communication and responsibility in educational activities.}
      \end{cvitems}
    }
\end{cventries}

\end{document}
